\clearpage
\section{Kesimpulan dan Saran}

\subsection{Kesimpulan}

Pada Tugas Besar 2 IF3170 Inteligensi Artifisial, tim berhasil membangun \textit{pipeline} lengkap untuk
memprediksi pembatalan pemesanan penginapan, mulai dari analisis data,
\textit{preprocessing}, pelatihan, hingga evaluasi. Lima algoritma, yaitu KNN,
\textit{Decision Tree Learning} (C4.5), \textit{Logistic Regression}, SVM, dan
\textit{Naive Bayes}, diimplementasikan dari awal menggunakan NumPy dan
dibandingkan dengan scikit-learn. Analisis data menunjukkan bahwa sekitar 33\% pemesanan
dibatalkan dan sekitar 39\% baris memiliki nilai kosong. Oleh karena itu
\textit{preprocessing} disesuaikan dengan kebutuhan tiap algoritma, dan
pembagian data memakai StratifiedGroupKFold agar baris yang identik
tidak berada pada data latih dan validasi sekaligus. Hasil implementasi sendiri
sesuai dengan scikit-learn. KNN dan \textit{Naive Bayes} menghasilkan
probabilitas yang hampir identik, dan selisih metrik \textit{Logistic Regression}
kurang dari 0,002. C4.5 berbeda karena algoritma pembandingnya adalah CART,
tetapi hasilnya tetap sebanding.

Berdasarkan \textit{5-fold cross-validation}, KNN (F1 macro 0,832) dan C4.5
(0,830) memberikan performa terbaik dengan selisih yang lebih kecil daripada
variasi antar-\textit{fold}, sehingga keduanya dapat dianggap setara.
Keduanya unggul karena mampu menangkap pola nonlinear dan interaksi antarfitur,
sedangkan \textit{Logistic Regression} (0,759), SVM berkernel linear (0,754),
dan \textit{Naive Bayes} (0,748) hanya membentuk batas yang lebih sederhana atau
mengasumsikan fitur saling bebas. Menambahkan pohon ketergantungan antarfitur
lewat \textit{Tree Augmented Naive Bayes} (bonus) menaikkan F1 macro
\textit{Naive Bayes} menjadi 0,775, sehingga asumsi independensi memang ikut
menurunkan skornya, tetapi belum cukup untuk menyamai KNN dan C4.5. Seluruh model berada jauh di atas
\textit{baseline} 0,675, tetapi \textit{recall} kelas batal baru mencapai 0,60
sampai 0,74. Pada bonus, penggeseran \textit{threshold} menjadi 0,667 untuk
\textit{cost-sensitive classification} menurunkan total biaya pada semua model,
dengan biaya terendah dicapai KNN (5.093). Visualisasi pohon, video lintasan
parameter \textit{Logistic Regression}, dan graf \textit{Naive Bayes} juga
berhasil dibuat untuk membantu menjelaskan perilaku model.

\subsection{Saran}

Untuk pengembangan lebih lanjut, disarankan mencoba model \textit{ensemble}
seperti \textit{Random Forest} atau \textit{gradient boosting} yang menggabungkan
kelebihan pohon keputusan, serta melakukan \textit{hyperparameter tuning} yang
lebih ekstensif. Rekayasa fitur tambahan, misalnya fitur interaksi untuk model
linear, juga berpotensi meningkatkan \textit{Logistic Regression} dan SVM.
Karena \textit{recall} kelas batal masih rendah, pemilihan \textit{threshold}
yang dioptimalkan langsung terhadap \textit{Macro F1} atau terhadap biaya
dapat dicoba. Pada \textit{Naive Bayes}, TAN menaikkan skor dengan memodelkan
ketergantungan antarfitur, sehingga model yang memodelkan interaksi lebih dari
dua fitur dapat dicoba untuk menyempitkan selisih dengan KNN dan C4.5.

Dari sisi implementasi, SVM masih dilatih pada 1.000 baris karena matriks
kernel penuh membutuhkan memori $O(n^2)$ sehingga perlu teknik seperti
\textit{kernel cache} dan \textit{shrinking} agar dapat memakai seluruh data.
Implementasi \textit{Logistic Regression} dengan \textit{Elastic Net} juga jauh
lebih lambat daripada scikit-learn sehingga dapat dipercepat dengan
laju belajar adaptif atau \textit{solver} yang lebih efisien. Terakhir, laporan
ini belum memuat skor \textit{leaderboard} Kaggle untuk seluruh model, sehingga
pengujian pada data uji perlu dilengkapi untuk memastikan hasil validasi dapat
digeneralisasi.
